\chapter{Lattice Geometry and Combinatorial Invariants}
\label{chap:lattice-architecture}

\section{Circular Lattices and Dihedral Symmetries}

Let $\mathbb{Z}_n = \{0, 1, \dots, n-1\}$ denote the cyclic ring of integers modulo $n$, geometrically embedded as the vertices of a regular $n$-gon inscribed in the unit circle and oriented clockwise with vertex $0$ positioned at the zenith. 

The physical isometries of the polygon constitute the dihedral group $D_n$. Geometrically, $D_n$ is generated by two reflections $\tau$ and $\tau' = \tau \rho$ whose reflection axes intersect at angle $\pi / n$. In the classical reflection (Coxeter) presentation, the group is given by:
\begin{equation}
D_n = \langle \rho, \tau \mid \rho^n = \tau^2 = (\tau \rho)^2 = 1 \rangle
\end{equation}
where $\rho = \tau \tau'$ is the clockwise unit rotation $x \mapsto x + 1 \pmod n$ and $\tau$ is the vertical chord reflection $x \mapsto -x \equiv n - x \pmod n$. In academic literature, the involution relation $(\tau \rho)^2 = 1$ is frequently rewritten as the conjugation relation:
\begin{equation}
\tau \rho \tau = \rho^{-1}
\end{equation}
emphasising the semidirect product decomposition $D_n \cong \mathbb{Z}_n \rtimes \mathbb{Z}_2$, in which the reflection $\tau \in \mathbb{Z}_2$ acts on the normal cyclic subgroup $\langle \rho \rangle \cong \mathbb{Z}_n$ by group inversion ($x \mapsto -x$). For the degenerate small rings, the isometry groups coincide with the physical symmetries:
\begin{equation}
D_1 = \{1\} \; (|D_1| = 1), \qquad D_2 = \langle \rho \mid \rho^2 = 1 \rangle \cong C_2 \; (|D_2| = 2)
\end{equation}
where on two vertices the single reflection $x \mapsto 1 - x \pmod 2$ acts identically to the transposition rotation $\rho$. For any $n \ge 1$, the reflections in $D_n$ are precisely the involutions $x \mapsto b - x \pmod n$ for $b \in \mathbb{Z}_n$.

A \textbf{composition mosaic} $\mathcal{M}$ on $\mathbb{Z}_n$ is a set partition of the $n$ vertices into $s$ non-empty, pairwise disjoint blocks:
\begin{equation}
\mathcal{M} = \{B_1, B_2, \dots, B_s\}, \quad \bigcup_{j=1}^s B_j = \mathbb{Z}_n, \quad B_i \cap B_j = \emptyset \; (i \neq j)
\end{equation}
Each constituent block $B_j \subset \mathbb{Z}_n$ of cardinality $d_j = |B_j| \ge 1$ forms an inscribed polygon in the lattice:
\begin{itemize}
    \item A \textbf{unigon} ($d_j = 1$), visualised as a central circular pip on vertex $v \in \mathbb{Z}_n$.
    \item A \textbf{digon} ($d_j = 2$), visualised as a straight metric chord connecting the vertex pair.
    \item An \textbf{inscribed polygon} ($d_j \ge 3$), visualised as a shaded convex polygonal tile bounded by circular perimeter chords.
\end{itemize}

\section{The Graded Stratification of Set Partitions}

Across the full combinatorial universe, the total number of set partitions of an $n$-element set is given by the $n$-th Bell number $B_n$:
\begin{equation}
B_n = \sum_{s=1}^n \genfrac{\{}{\}}{0pt}{}{n}{s}
\end{equation}
where $\genfrac{\{}{\}}{0pt}{}{n}{s}$ denotes the Stirling numbers of the second kind, enumerating set partitions into exactly $s$ blocks. 

The space of partitions further stratifies according to the unordered multiset of block sizes, termed the \textbf{partition shape} $\lambda = (\lambda_1, \dots, \lambda_s) \vdash n$, where $\lambda_1 \ge \lambda_2 \ge \dots \ge \lambda_s \ge 1$ and $\sum_{j=1}^s \lambda_j = n$. For each shape $\lambda \vdash n$, the exact number of set partitions is governed by the Fa\`a di Bruno multinomial formula:
\begin{equation}
p_\lambda = \frac{n!}{\prod_{i=1}^n (i!)^{c_i} \, c_i!}
\end{equation}
where $c_i = c_i(\lambda)$ is the multiplicity of parts of size $i$ in $\lambda$.

Across the first nine rings $\mathbb{Z}_1$ through $\mathbb{Z}_9$, the Bell numbers grow rapidly:
\begin{gather*}
B_1 = 1, \quad B_2 = 2, \quad B_3 = 5, \quad B_4 = 15, \quad B_5 = 52 \\
B_6 = 203, \quad B_7 = 877, \quad B_8 = 4{,}140, \quad B_9 = 21{,}147
\end{gather*}
yielding a core combinatorial universe of $26{,}442$ set partitions across $96$ integer partition shapes. Extending this progression through $\mathbb{Z}_{16}$ reveals the full exponential escalation of the partition continuum, reaching $B_{16} = 10{,}480{,}142{,}147$ and encompassing a cumulative universe of $12{,}086{,}679{,}035$ set partitions across $914$ partition shapes.

\section{The Dihedral Orbit Inventory and Burnsidean Enumeration}

The dihedral group $D_n$ acts on the space of set partitions $\mathcal{P}(\mathbb{Z}_n)$ by permuting the underlying coordinates:
\begin{equation}
g \cdot \{B_1, \dots, B_s\} = \{g(B_1), \dots, g(B_s)\} \quad (g \in D_n)
\end{equation}
Because the group action is an isometry, it preserves the part count $s$, the individual block cardinalities $|B_j|$, and the partition shape $\lambda \vdash n$. The set partitions of shape $\lambda$ decompose into disjoint $D_n$-orbits:
\begin{equation}
\mathcal{O}(\mathcal{M}) = \{g \cdot \mathcal{M} \mid g \in D_n\}
\end{equation}

By the Orbit-Stabiliser Theorem, the length of any orbit $|\mathcal{O}| = k$ must divide the group order $|D_n|$. The stabiliser subgroup $\operatorname{Stab}_{D_n}(\mathcal{M})$ has order $|D_n| / k$.

Summing the orbit multiplicities $a_{k,\lambda}$ yields the graded bivariate orbit polynomial $P_n(\mathbf{z}, \omega)$:
\begin{equation}
P_n(\mathbf{z}, \omega) = \sum_{\lambda \vdash n} Q_\lambda(\omega) \mathbf{z}^\lambda = \sum_{\lambda \vdash n} \left( \sum_{k \mid |D_n|} a_{k,\lambda} \omega^k \right) \mathbf{z}^\lambda
\end{equation}
Evaluating this polynomial at $\omega = 1$ recovers the total count $M_n$ of non-equivalent dihedral composition mosaics on $\mathbb{Z}_n$. Table~\ref{tab:census-z1-z16} presents the consolidated inventory across the sixteen rings $\mathbb{Z}_1$ through $\mathbb{Z}_{16}$, recovering the initial sixteen terms of OEIS sequence A084708 and establishing the grand total of $\mathbf{382{,}030{,}946}$ non-equivalent dihedral composition mosaics (of which the initial nine rings comprise the $1{,}859$ visually catalogued mosaics of the core atlas).

% Consolidated Dihedral Mosaic Census from Z1 to Z16
\begin{table}[htbp]
\centering
\small
\setlength{\tabcolsep}{4.5pt}
\caption{\label{tab:census-z1-z16}\label{tab:census-z1-z8}Consolidated Dihedral Mosaic Census from $\mathbb{Z}_1$ to $\mathbb{Z}_{16}$ (OEIS A084708)}
\begin{tabular}{c r r r r l}
\toprule
Ring $\mathbb{Z}_n$ & $|D_n|$ & Partitions $B_n$ & Shapes & Dihedral Mosaics & Orbit Size Decomposition $\sum a_k \omega^k$ \\
\midrule
$\mathbb{Z}_1$ & 1 & 1 & 1 & \textbf{1} & $1\omega^{1}$ \\
$\mathbb{Z}_2$ & 2 & 2 & 2 & \textbf{2} & $2\omega^{1}$ \\
$\mathbb{Z}_3$ & 6 & 5 & 3 & \textbf{3} & $2\omega^{1} + 1\omega^{3}$ \\
$\mathbb{Z}_4$ & 8 & 15 & 5 & \textbf{7} & $3\omega^{1} + 2\omega^{2} + 2\omega^{4}$ \\
$\mathbb{Z}_5$ & 10 & 52 & 7 & \textbf{12} & $2\omega^{1} + 10\omega^{5}$ \\
$\mathbb{Z}_6$ & 12 & 203 & 11 & \textbf{37} & \begin{tabular}[t]{@{}l@{}} $4\omega^{1} + 2\omega^{2} + 9\omega^{3} + {}$ \\ $\quad 16\omega^{6} + 6\omega^{12}$ \end{tabular} \\
$\mathbb{Z}_7$ & 14 & 877 & 15 & \textbf{93} & $2\omega^{1} + 57\omega^{7} + 34\omega^{14}$ \\
$\mathbb{Z}_8$ & 16 & 4{,}140 & 22 & \textbf{354} & \begin{tabular}[t]{@{}l@{}} $4\omega^{1} + 6\omega^{2} + 27\omega^{4} + {}$ \\ $\quad 132\omega^{8} + 185\omega^{16}$ \end{tabular} \\
$\mathbb{Z}_9$ & 18 & 21{,}147 & 30 & \textbf{1{,}350} & \begin{tabular}[t]{@{}l@{}} $3\omega^{1} + 13\omega^{3} + 323\omega^{9} + {}$ \\ $\quad 1{,}011\omega^{18}$ \end{tabular} \\
$\mathbb{Z}_{10}$ & 20 & 115{,}975 & 42 & \textbf{6{,}351} & \begin{tabular}[t]{@{}l@{}} $4\omega^{1} + 3\omega^{2} + 103\omega^{5} + {}$ \\ $\quad 937\omega^{10} + 5{,}304\omega^{20}$ \end{tabular} \\
$\mathbb{Z}_{11}$ & 22 & 678{,}570 & 56 & \textbf{31{,}950} & \begin{tabular}[t]{@{}l@{}} $2\omega^{1} + 2{,}208\omega^{11} + {}$ \\ $\quad 29{,}740\omega^{22}$ \end{tabular} \\
$\mathbb{Z}_{12}$ & 24 & 4{,}213{,}597 & 77 & \textbf{179{,}307} & \begin{tabular}[t]{@{}l@{}} $6\omega^{1} + 11\omega^{2} + 31\omega^{3} + 38\omega^{4} + {}$ \\ $\quad 288\omega^{6} + 11\omega^{8} + {}$ \\ $\quad 6{,}885\omega^{12} + 172{,}037\omega^{24}$ \end{tabular} \\
$\mathbb{Z}_{13}$ & 26 & 27{,}644{,}437 & 101 & \textbf{1{,}071{,}265} & \begin{tabular}[t]{@{}l@{}} $2\omega^{1} + 16{,}031\omega^{13} + {}$ \\ $\quad 1{,}055{,}232\omega^{26}$ \end{tabular} \\
$\mathbb{Z}_{14}$ & 28 & 190{,}899{,}322 & 135 & \textbf{6{,}845{,}581} & \begin{tabular}[t]{@{}l@{}} $4\omega^{1} + 4\omega^{2} + 1{,}042\omega^{7} + {}$ \\ $\quad 53{,}918\omega^{14} + 6{,}790{,}613\omega^{28}$ \end{tabular} \\
$\mathbb{Z}_{15}$ & 30 & 1{,}382{,}958{,}545 & 176 & \textbf{46{,}162{,}583} & \begin{tabular}[t]{@{}l@{}} $4\omega^{1} + 18\omega^{3} + 162\omega^{5} + 2\omega^{6} + {}$ \\ $\quad 118\omega^{10} + 127{,}459\omega^{15} + {}$ \\ $\quad 46{,}034{,}820\omega^{30}$ \end{tabular} \\
$\mathbb{Z}_{16}$ & 32 & 10{,}480{,}142{,}147 & 231 & \textbf{327{,}731{,}950} & \begin{tabular}[t]{@{}l@{}} $5\omega^{1} + 13\omega^{2} + 121\omega^{4} + 3{,}422\omega^{8} + {}$ \\ $\quad 449{,}637\omega^{16} + 327{,}278{,}752\omega^{32}$ \end{tabular} \\
\bottomrule
\end{tabular}
\end{table}


\section{The Bilateral Symmetry Classification}

Every individual block $B \subset \mathbb{Z}_n$ within a mosaic possesses an internal circular stride composition $(s_1, \dots, s_d)$ satisfying $\sum_{i=1}^d s_i = n$. A block is defined as \textbf{bilaterally symmetric} (or achiral) if it is invariant under at least one reflection $x \mapsto b - x \pmod n$ of the circular lattice. 

Under this classification:
\begin{itemize}
    \item \textbf{Achiral Blocks (Blue)}: Inscribed polygons possessing an internal reflection axis. On any lattice $\mathbb{Z}_n$, every unigon ($d=1$), every digon chord ($d=2$), and every regular or isosceles polygon is strictly achiral.
    \item \textbf{Chiral Blocks (Pink)}: Scalene polygons whose circular stride sequence is not invariant under cyclic reversal.
\end{itemize}

\begin{theorem}[Emergence of Block Chirality]
On the circular lattices $\mathbb{Z}_n$, all constituent blocks of all set partitions are strictly achiral for $n \le 5$. Chiral blocks emerge if and only if $n \ge 6$ and block cardinality $d \ge 3$. The minimal chiral block is the scalene triangle of circular strides $(1, 2, 3)$ on the hexagonal lattice $\mathbb{Z}_6$.
\end{theorem}

\begin{proof}
For block size $d \in \{1, 2\}$, any unigon is trivially self-symmetric, and any digon $\{u, v\}$ is invariant under the reflection through the perpendicular bisector of the chord $u \leftrightarrow v$, which is a valid lattice reflection for every $n$. 

For $d = 3$, an inscribed triangle with stride composition $(s_1, s_2, s_3)$ is chiral if and only if the strides are pairwise distinct (scalene). The minimal sum of three distinct positive integers is $1 + 2 + 3 = 6$. Hence no chiral triangle can exist for $n < 6$. For $n = 6$, the unique integer partition into three distinct parts is $1 + 2 + 3 = 6$, realised by the block $\{0, 1, 3\}$ with strides $(1, 2, 3)$. Under reflection $x \mapsto b - x \pmod 6$, the stride sequence reverses to $(3, 2, 1) \equiv (1, 3, 2) \not\equiv (1, 2, 3)$ under cyclic shifts. Hence $\{0, 1, 3\}$ possesses no internal reflection symmetry and is chiral.
\end{proof}

A mosaic $\mathcal{M}$ is characterised by its \textbf{chirality signature}, recording the binary achiral/chiral character of its sorted constituent blocks (e.g. $[ssss]$ for an all-symmetric mosaic, $[sscc]$ for a mosaic containing two symmetric and two chiral blocks). In the visual atlas, chiral blocks appear in distinctive rose-pink shading, highlighting the subtle symmetry-breaking phenomena that first emerge on $\mathbb{Z}_6$.

\section{The Singleton Shattering Duality}

A foundational structural property uniting all cyclic lattices is the \textbf{Singleton Shattering Duality} (in the spirit of Joyal's combinatorial species and Kerber's permutation representations). For any partition shape $\lambda = (b_1, \dots, b_{s-1}, a) \vdash n$ possessing a distinguished block of size $a \ge 2$, shattering this block into $a$ individual singleton vertices ($a \to 1^a$) forms an equivariant bijection between point stabilisers.

Consequently, the orbit polynomial is identically preserved under singleton replacement:
\begin{equation}
Q_{(b_1, \dots, b_{s-1}, a)}(\omega) \equiv Q_{(b_1, \dots, b_{s-1}, 1^a)}(\omega)
\end{equation}

In particular, for any two-block partition shape $n = a + b$ with distinct block cardinalities $a \ne b$, complementation on the Boolean lattice of subsets induces a \textbf{tripartite shattering triad}:
\begin{equation}
z_a z_b \;\longleftrightarrow\; z_a z_1^b \;\longleftrightarrow\; z_b z_1^a
\end{equation}
Because any lattice isometry $g \in D_n$ fixes a block $A$ setwise if and only if it fixes its complement $B = A^c$, and because $a \ne b$ prevents $g$ from interchanging $A$ and $B$, the stabiliser subgroups satisfy:
\begin{equation}
\operatorname{Stab}_{D_n}(\{A, B\}) = \operatorname{Stab}_{D_n}(A) = \operatorname{Stab}_{D_n}(B) = \operatorname{Stab}_{D_n}(\{A\} \cup \{\{v\} \mid v \in B\})
\end{equation}
Hence all three partition shapes share identical orbit polynomials $Q(\omega)$ across all cyclic lattices:
\begin{align}
\text{On } \mathbb{Z}_4\colon \quad z_4 &\longleftrightarrow z_1^4 && (1\omega^1) \\
\text{On } \mathbb{Z}_5\colon \quad z_5 &\longleftrightarrow z_1^5 && (1\omega^1) \\
\text{On } \mathbb{Z}_6\colon \quad z_6 &\longleftrightarrow z_1^6 && (1\omega^1) \\
z_2 z_4 &\longleftrightarrow z_2 z_1^4 \longleftrightarrow z_4 z_1^2 && (1\omega^3 + 2\omega^6) \\
\text{On } \mathbb{Z}_7\colon \quad z_2 z_5 &\longleftrightarrow z_2 z_1^5 \longleftrightarrow z_5 z_1^2 && (3\omega^7) \\
z_3 z_4 &\longleftrightarrow z_3 z_1^4 \longleftrightarrow z_4 z_1^3 && (3\omega^7 + 1\omega^{14}) \\
\text{On } \mathbb{Z}_8\colon \quad z_8 &\longleftrightarrow z_1^8 && (1\omega^1) \\
z_2 z_6 &\longleftrightarrow z_2 z_1^6 \longleftrightarrow z_6 z_1^2 && (1\omega^4 + 3\omega^8) \\
z_3 z_5 &\longleftrightarrow z_3 z_1^5 \longleftrightarrow z_5 z_1^3 && (3\omega^8 + 2\omega^{16}) \\
z_2^2 z_4 &\longleftrightarrow z_2^2 z_1^4 && (1\omega^2 + 4\omega^4 + 10\omega^8 + 7\omega^{16})
\end{align}

\section{Circular Chord Energy}

The spatial dispersion of a composition mosaic $\mathcal{M} = \{B_1, \dots, B_s\}$ on $\mathbb{Z}_n$ is quantified by its \textbf{chord dispersion energy} $E(\mathcal{M})$:
\begin{equation}
E(\mathcal{M}) = \sum_{B \in \mathcal{M}} \sum_{u < v \in B} \min(|u - v|, n - |u - v|)^2
\end{equation}
defined as the sum of squared circular shortest-path distances across all internal edges of every constituent block. 

Configurations whose vertices are closely clustered span small perimeter strides and exhibit low energy, whereas configurations whose vertices are diametrically distributed across the lattice maximise chord dispersion. Within any given partition shape, chord energy provides a natural geometric gradient for canonical ordering.
